% Two-column conference paper. Swap \documentclass for the venue's own class when you have it:
% the structure below survives the change.
\documentclass[10pt,twocolumn,letterpaper]{article}
\usepackage[margin=0.75in]{geometry}
\usepackage{graphicx}
\usepackage{amsmath,amssymb}
\usepackage{booktabs}
\usepackage[hidelinks]{hyperref}
\usepackage[capitalise]{cleveref}

\title{A title that states the result}
\author{First Author\\ Affiliation\\ \texttt{first@example.edu}
  \and Second Author\\ Affiliation\\ \texttt{second@example.edu}}
\date{}

\begin{document}
\maketitle

\begin{abstract}
One sentence on the problem, one on what is new, one on the evidence, one on why it matters.
\end{abstract}

\section{Introduction}
\label{sec:intro}

\section{Related work}
\label{sec:related}

\section{Method}
\label{sec:method}

\begin{equation}
  \label{eq:objective}
  \mathcal{L}(\theta) = \mathbb{E}_{x \sim \mathcal{D}} \left[ \ell(f_\theta(x), y) \right].
\end{equation}

\section{Experiments}
\label{sec:experiments}

\begin{table}[t]
  \centering
  \caption{Results. Bold is best.}
  \label{tab:main}
  \begin{tabular}{lrr}
    \toprule
    Method & Accuracy & Runtime (s) \\
    \midrule
    Baseline & 00.0 & 0.0 \\
    Ours     & \textbf{00.0} & 0.0 \\
    \bottomrule
  \end{tabular}
\end{table}

\Cref{tab:main} reports the main comparison; \cref{eq:objective} is the objective it optimises.

\section{Conclusion}
\label{sec:conclusion}

\bibliographystyle{plain}
\bibliography{refs}
\end{document}
